%%%% RESUMO
%%
%% Apresentação concisa dos pontos relevantes de um texto, fornecendo uma visão rápida e clara do conteúdo e das conclusões do
%% trabalho.

\begin{resumoutfpr}%% Ambiente resumoutfpr

A fusão de imagens para construção de mosaicos panorâmicos é uma técnica fundamental em visão computacional com aplicações em vigilância, navegação autônoma e monitoramento ambiental. Apesar dos avanços recentes com técnicas de aprendizado profundo, falta na literatura uma comparação sistemática entre abordagens clássicas e modernas considerando múltiplos critérios de qualidade e eficiência. Além disso, o impacto da qualidade de fusão em tarefas subsequentes como detecção de objetos permanece pouco explorado. Este trabalho apresenta uma análise comparativa abrangente de 20 técnicas de fusão de imagens, sendo 15 baseadas em métodos clássicos e 5 em aprendizado profundo, aplicadas ao dataset Robust Image Stitching with Multiple Registrations. Foram realizadas 295 execuções bem-sucedidas sobre 15 conjuntos de imagens, com avaliação quantitativa através de métricas PSNR, SSIM, MSE, Ghosting Score, Seam Energy, DeltaE, Entropia, Nitidez e tempo de processamento. Os resultados mostram que os métodos OpenCV Stitcher e Stitcher Básico alcançaram melhor desempenho geral (score 0,8732), com PSNR de 11,10 dB e tempo médio de 0,26 segundos, representando excelente equilíbrio entre qualidade e eficiência. O método Seams Exposure Multiband apresentou o maior PSNR (11,49 dB) e menor energia de seam (77,42), indicando melhor qualidade visual, porém com custo computacional significativamente maior (7,11s). A comparação entre abordagens clássicas e de aprendizado profundo revelou diferença pequena em qualidade (+0,08 dB PSNR, +0,0066 SSIM para deep learning), mas com tempo de processamento médio 68\% maior para as técnicas baseadas em redes neurais. Todos os métodos apresentaram Ghosting Score alto (~0,94), refletindo os desafios inerentes ao dataset utilizado. Na segunda parte do trabalho, investiga-se o impacto das diferentes técnicas de fusão na detecção de objetos, aplicando modelos YOLO (YOLOv5, YOLOv8 e YOLO-NAS) sobre os panoramas gerados e comparando com detecções em imagens individuais. \textcolor{blue}{[Resultados de detecção serão adicionados após conclusão dos experimentos.]} Os resultados demonstram que técnicas clássicas oferecem excelente relação custo-benefício para fusão panorâmica, sendo adequadas para aplicações que requerem processamento rápido, enquanto métodos de aprendizado profundo são indicados quando qualidade visual é prioridade absoluta. O trabalho fornece diretrizes práticas para escolha de técnicas de fusão considerando o contexto de aplicação e estabelece uma base metodológica para investigações futuras sobre integração de fusão de imagens e detecção de objetos.

\vspace{1em}

\noindent\textbf{Palavras-chave}: fusão de imagens; mosaico panorâmico; image stitching; visão computacional; aprendizado profundo; detecção de objetos; YOLO.

\end{resumoutfpr}
